\documentclass[12pt]{article}
\usepackage{xeCJK}
\usepackage{fontspec}
\usepackage{amsmath}
\usepackage{graphicx}
\usepackage{geometry}
\usepackage{siunitx}
\usepackage{float}
\usepackage{fancyhdr}
\usepackage{circuitikz}
\usepackage{caption}
\usepackage{subcaption}
\usepackage{lastpage}
\usepackage{tikz}
\captionsetup[figure]{
    font=footnotesize,
    name=图
}
\captionsetup[table]{
    font=footnotesize,
    name=表
}
\setCJKmainfont{Songti SC}
\setCJKfamilyfont{kaishu}{Kaiti SC} % 定义楷书字体族
\geometry{a4paper, margin=1in, top=1in}
\setlength{\parindent}{12pt}
\setlength{\headheight}{13.6pt}
% 页眉页脚设置
\pagestyle{fancy}
\fancyhf{}
\fancyhead[R]{\small 普通物理实验报告}
\fancyhead[C]{\small 实验名称：弗兰克-赫兹实验}
\fancyfoot[C]{\small 第 \thepage\ 页/共 \pageref{LastPage} 页}
\fancypagestyle{plain}{%
    \fancyhf{}
    \fancyfoot[C]{\small 第 \thepage\ 页/共 \pageref{LastPage} 页}
    \renewcommand{\headrulewidth}{0pt}  % 移除页眉分割线
}
\newcommand{\question}[1]{{\CJKfamily{kaishu}#1}}
\title{\Large\textbf {弗兰克-赫兹实验自主学习报告}}
\author{普通物理实验2\\
        姓名：\underline{\hspace{1cm}[REDACTED]\hspace{1cm}} \\
        学号：\underline{\hspace{0.5cm}[REDACTED]\hspace{0.5cm}} \\
        序号：\underline{\hspace{1.3cm}[REDACTED]\hspace{1.3cm}}}
\date{\today}
\begin{document}
\maketitle
\begin{enumerate}
    \item \question{弗兰克-赫兹实验的主要实验现象是什么？该现象说明了什么物理问题？}
    \par
    \hspace{1em}主要实验现象是，阳极电流 $I_A$ 随加速电压 $U_{Kg2}$ 的增加并不是单调地上升，而是呈现出周期性的升降现象，且各极大值之间的间距相同。
    \par
    \hspace{1em}具体表现为当电压达到特定值（汞原子的第一激发电位）时，电流会突然显著下降 。如果继续增大加速电压，电流会再次上升，并在电压达到第一激发电位的整数倍附近时，再次出现急剧下降的现象。\\
    \par
    \hspace{1em}该现象说明了原子内部的能量状态是不连续的，证明了原子量子化能级的存在，即玻尔能级。
    \par
    \hspace{1em}当电子动能小于第一激发电位时，电子与汞原子发生的是弹性碰撞，能量交换极小；当电子动能达到第一激发电位时，发生非弹性碰撞，汞原子跃迁到激发态，电子自身能量减小，无法到达阳极，从而引起电流下降。

    \item \question{说明四栅式弗兰克-赫兹管的基本结构和各个电极的作用:}
    \par
   \hspace{1em}四栅式F-H管的基本结构包括（灯丝）、阴极K、第一栅极g1、第二栅极g2和阳极A 。
    \begin{itemize}
        \item \textbf{阴极 K }：经灯丝F加热后通过热电子发射效应发射电子 。
        \item \textbf{第一栅极 g1 }：g1电位略高于K，主要作用是消除阴极附近因电子云堆积形成的空间电荷效应（电子堆积）对电子发射的抑制，使电子顺利从阴极发射。
        \item \textbf{第二栅极 g2 }：是主要的加速电极。在其与阴极之间施加一个可变的加速电压 $U_{Kg2}$，电子在此区域被加速获得并不断与汞原子碰撞 。
        \item \textbf{阳极 p }：收集电子。g2与p之间存在一个减速电压 。其作用是只允许那些动能足够大的电子（即未在g2附近发生非弹性碰撞的电子）能够克服反向电场到达阳极形成电流，这使得电流下降的现象更加显著，提升实验图像的特征衬比。
    \end{itemize}

    \item \question{实验中如何做到只观察汞的第一激发态？电路和参数设置上有什么考虑？如要观察高激发态，电路和参数设置上要做什么调整？}
    \par
    \hspace{1em}在F-H标准实验中，我们一般只能观察到汞的第一激发态引起的现象，原因是此时汞管温度高，汞蒸汽压大，电子的平均自由程较小，在达到足以汞高激发态的动能之前一般会发射碰撞能量减小，电子动能很难超过4.9eV。以下参数设置较为重要：
    \begin{itemize}
        \item \textbf{灯丝电压 $U_F$ 和炉温}：灯丝电压$3 \sim 5V$，炉温控制在180℃左右，以保证适宜的汞蒸气密度和电子密度。
        \item \textbf{第一栅极电压 $U_{Kg1}$}：不宜过小，以消除空间电荷效应。
        \item \textbf{减速电压 $U_{g2p}$}：设置为1V左右（实验中调节），以有效筛选出发生非弹性碰撞的电子，使现象更明显的值即可。
        \item \textbf{电路设置}：使用微弱电流放大器来放大阳极收集到的电流信号。

    \end{itemize}
    \par
    \hspace{1em}如要观察高激发态，难度较大，因为电子能量达到4.9eV后，激发第一能级的几率远大于激发高能级的几率。通常需要对实验方法进行以下调整：
    \begin{itemize}
        \item \textbf{炉温设置：} 炉温足够低（130℃左右），减小汞原子的密度，增大电子的自由程使得电子在碰撞前获得较高的动能。
        \item \textbf{电路设置：} 在K和g1间加加速电压（距离很短），这样就可以减少电子达到高动能前的碰撞，g1,g2电压差很小，电子在两者间的空间与汞原子发生碰撞，g2p仍为减速电压。
        \item \textbf{其他设置：} 加速电压范围更大，几十伏特;换用更高灵敏度的微电流放大器
    \end{itemize}

    
    \item \question{微电流放大器为什么可以同时获得高放大倍数和小输入阻抗？}
    \par
    微电流放大器使用了运算放大器的\textbf{负反馈}工作模式。
    \par
    利用运放的“虚地”特性，输入端电位$\approx 0V$，对于信号源来说，运算放大器阻抗非常小，实现了信号和后端的隔离。
    再利用非常大的反馈电阻$Rf$,有关系式$\frac{-U}{R_f}+I_i = I$又$I_iR_i=U_{-}\approx 0$，故$｜U｜=IR_f$非常大，实现了高放大倍数。
    
    \item \question{（选做）如果希望测量汞的电离能，电路和参数应当如何设置？}
    \par
    \hspace{1em}电子电离电位10.4eV比第一激发电位4.9eV大很多。电子达到10.4eV与汞原子碰撞后会产生电子雪崩效应，阳极电流 $I_A$ 出现急剧增长。
    \par
    \hspace{1em}同样，为了避免电子提前碰撞损失能量，电路设置应该同测量高激发态能级的配置，电子迅速加速到指定能量。其他需要的操作有。
    \begin{itemize}
        \item 将加速电压 $U_{Kg}$ 的扫描范围扩大，使其能够超过10.4V，扫描至20V或更高。
        \item 观察 $I_A-U_{Kg}$ 曲线，寻找电流开始急剧上升的电压拐点。该拐点对应的电压值即为汞的电离电位。
    \end{itemize}
\end{enumerate}
\end{document}